\lecture{5}
\title{Context-Free Pumping Lemma, Turing Machines}

\begin{document}

Problem set 2 is due on Tuesday, October 13th at noon.
Quiz 1 is from 7:30 PM to 9 PM in Walker on October 6th.

\subsection{Context-Free Languages: Miscellaneous}

Recall from \href{/18.404/lectures/4/#cfls-are-pda-languages}{Lecture 4}
the following theorem (for which we only proved one direction).

\begin{theorem}{CFL $=$ PDA}
    A language is context-free if and only if
    it is recognized by some PDA.
\end{theorem}

The PDA $\implies$ CFL direction is proven by
a similar strategy as used to show GNFA $\implies$ Regular Expression.

It turns out that the following are also true;
these are optional exercises in problem set 2.
\begin{itemize}
    \item The class of CFLs is closed under union $A \cup B$,
    concatenation $AB$, and star $A^*$.
    \item The class of CFLs is \emph{not} closed
    under intersection $A \cap B$ and complementation $\overline{A}$.
    \begin{itemize}
        \item Roughly, $A \cap B$ is not necessarily a CFL
        because simulating the PDA of $A$ and the PDA of $B$
        requires \emph{two stacks}.
    \end{itemize}
    \item The intersection of a CFL $A$ and a regular language $L$
    is also a CFL.
    \begin{itemize}
        \item Roughly, $A \cap L$ is a CFL
        because simulating the PDA of $A$ and the NFA of $L$
        still only requires \emph{one stack}.
    \end{itemize}
\end{itemize}


\begin{remark}
    Outside of context-free languages,
    one could also talk about \emph{context-sensitive} languages
    whose rules look like:
\begin{verbatim}
      A ::= 'a' B 'b' | 'c'     // (1)
  A 'a' ::= B 'c'               // (2)
'c' A B ::= 'x'                 // (3)
\end{verbatim}
    Rules (2) and (3) depend on the surrounding context;
    these rules are not allowed for context-free languages.
\end{remark}

\subsection{Context-Free Pumping Lemma}

Just as we had a pumping lemma for regular languages,
we have one for CFLs too.

\begin{theorem}{CFL Pumping Lemma}
    For every CFL $A$,
    there must be some \textbf{pumping length} $p$ such that:
    \begin{quote}
        For all $s \in A$ with $|s| \geq p$,
        we may write $s = uvxyz$ with $vy \neq \epsilon$
        such that $|vxy| \leq p$ and
        $uv^ixy^iz \in A$ for all $i \in \mathbb{N}$.
    \end{quote}
\end{theorem}

\begin{proof}
    Any very long string $s \in A$
    must have a very tall parse tree.
    In particular, this very tall parse tree must have
    some variable $R$ that appears twice
    along some downward branch of the tree.
    \begin{center}
        \frame{\includegraphics[width=0.75\textwidth]{lecture-05/pumping-parse-tree}}
    \end{center}
    So long as such a variable $R$ is found,
    we can pump as pictured above;
    in short, duplicate the subtree that goes $R \to R$.
    There remain some details to be filled in---and, furthermore,
    some technical fixes our construction might need.

    \textbf{Question.}
    What should $p$ be such that
    (i) some variable $R$ must appear twice down a single path, and
    (ii) $|vxy| \leq p$ must hold?
    \begin{quote}
        Let $b$ denote the length of the longest replacement string
        in the CFG, and let $V$ be its variables.
        We claim $p = b^{|V| + 1}$ works.

        Consider the \emph{longest} root-to-leaf path in the parse tree.
        Since $|s| \geq p = b^{|V| + 1}$,
        this path must have height at least $|V| + 1$.
        So among the \emph{lowest} $|V| + 1$ edges of this path,
        some variable $R$ must appear twice, as needed.

        Furthermore, because the higher of the two $R$s
        has height \emph{at most} $|V| + 1$,
        it must be that $|vxy| \leq b^{|V| + 1} = p$. ~ $\square$
    \end{quote}

    \textbf{Question.}
    How can we guarantee that $vy \neq \epsilon$ in our construction?
    \begin{quote}
        Suppose that $v = \epsilon$ and $y = \epsilon$.
        Then the parse tree can be shrunk (without changing $s$)
        by replacing the subtree rooted at the upper $R$
        with the subtree rooted at the lower $R$.

        So we can prevent $vy = \epsilon$
        by only considering the parse tree of $s$
        with the fewest nodes. ~ $\square$
    \end{quote}

    And that's the proof. ~ $\blacksquare$
\end{proof}

\begin{problem}
    Prove that the language
    $C := \{\str{a}^k \str{b}^k \str{c}^k \mid k \geq 0\}$
    is not a CFL.
\end{problem}

\begin{solution}
    Suppose FTSOC that it is a CFL with pumping length $p$.
    Then consider $s = \str{a}^p\str{b}^p\str{c}^p \in C$.

    Say it splits like $s = uvxyz$ as in the pumping lemma.
    Since $|vxy| \leq p$, the substring $vxy$
    can cover at most two of the three blocks
    $\{\str{a}^p, \str{b}^p, \str{c}^p\}$.
    So pumping up must increase some letter's frequency
    while leaving another letter's frequency unchanged.
    But then the resulting word is not in $C$, contradiction. ~ $\blacksquare$
\end{solution}

\begin{remark}
    Earlier in this lecture, we said
    $A \cap B$ is not necessarily a CFL because
    it necessitates a PDA architecture with \emph{two stacks}.
    This example shows why \emph{two stacks} are stronger than
    \emph{one stack}---the language $C$ is easily recognizable
    given two stacks!
\end{remark}

\begin{problem}
    Prove that the language
    $D := \{w \mid w \text{ has equal \#s of }
    \str{a}, \str{b}, \str{c}\}$
    is not a CFL.
\end{problem}

\begin{solution}
    You \emph{could} use the same pumping-lemma contradiction
    $\str{a}^p\str{b}^p\str{c}^p$ as in the previous proof\dots
    or you could recall:
    \begin{quote}
        \textbf{Proposition.}
        The intersection of a CFL and a regular language is also a CFL.
    \end{quote}
    So if $D$ were a CFL, then so would be
    $D \cap \str{a}^*\str{b}^*\str{c}^* = C$.
    But we know $C$ is not a CFL, contradiction. ~ $\blacksquare$
\end{solution}

Here's an example where the choice of string $s$ is less obvious.

\begin{problem}
    Prove that the language
    $B := \{ww \mid w \in \{\str{0}, \str{1}\}^*\}$ is not a CFL.
\end{problem}

\begin{solution}
    You can check that
    $s = \str{0}^p\str{1}\str{0}^p\str{1} \in B$ does \emph{not} raise
    a contradiction via the pumping lemma.  Uh oh!

    But the string $s = \str{0}^p\str{1}^p\str{0}^p\str{1}^p \in B$
    does raise a contradiction, so we win anyway. ~ $\blacksquare$
\end{solution}

\subsection{Turing Machines}

A \textbf{Turing machine (TM)} is like a PDA, but instead of owning a \emph{stack},
the machine owns a \emph{read/write tape}.

\begin{center}
    \frame{\includegraphics[width=0.5\textwidth]{lecture-05/turing-machine}}
\end{center}

\textbf{Definition.}
A Turing machine is a 7-tuple
$M = (Q, \Sigma, \Gamma, \delta, q_0, q_{\mathrm{acc}}, q_{\mathrm{rej}})$,
where:
\begin{itemize}
    \item The usual NFA components have the same meanings as before\dots
    except for $q_{\mathrm{acc}}$ and $q_{\mathrm{rej}}$.
    \begin{itemize}
        \item The set of \textbf{states} is $Q$.
        \item The \textbf{input alphabet}
        (the set of atoms accepted as input) is $\Sigma$.
        \item The \textbf{start state} is $q_0 \in Q$.
        \item The \textbf{accept state} is $q_{\mathrm{acc}}$.  It is
        a \emph{halt state}, meaning the machine halts
        upon reaching $q_{\mathrm{acc}}$.
        \item The \textbf{reject state} is $q_{\mathrm{rej}}$.  It is
        a \emph{halt state}, meaning the machine halts
        upon reaching $q_{\mathrm{rej}}$.
    \end{itemize}
    \item The \textbf{tape alphabet}
    (the set of atoms accepted by the tape) is $\Gamma$,
    which must contain a \textbf{blank symbol} $\sqcup$
    and everything in $\Sigma$.
    \begin{itemize}
        \item The tape initially holds the input word $w$,
        followed by infinitely many instances of the blank symbol $\sqcup$.
    \end{itemize}
    \item The \textbf{transition function} has type
    $\delta \colon Q \times \Gamma \to
    Q \times \Gamma \times \{\mathrm{L}, \mathrm{R}\}$.
    So at each step, the Turing machine\dots
    \begin{itemize}
        \item \dots reads in the current state $\in Q$
        and the atom $\in \Gamma$ at its tape position \dots
        \item \dots then moves to a new state $\in Q$,
        writes an atom $\in \Gamma$ at its tape position,
        and moves left or right along the tape.
    \end{itemize}
\end{itemize}
\emph{\textcolor{red}{Warning!}} ~ 
Whereas PDAs and FAs read in the input
one character at a time from left to right\dots
Turing machines read in the input off of the tape itself,
and may read it in any way they choose.
(So in the example image above,
the input $w$ is \str{babaisMe}.)

Note that for any word $w$, a Turing machine $M$ either
\textbf{accepts} $w$, \textbf{rejects} $w$, or \textbf{never halts}.

\begin{remark}
    There is no sense in which a Turing machine
    ``finishes reading the input''---again,
    because Turing machines do not read the input from left-to-right
    like PDAs and FAs do.
    So halting can \emph{only} occur
    by arriving at the accept or reject state.
\end{remark}

\textbf{Definition.}
A language $A$ is \textbf{Turing-recognizable} if
$A = L(M)$ for some Turing machine $M$.

\textbf{Definition.}
A language $A$ is \textbf{decidable} if
$A = L(M)$ for some Turing machine $M$ that \emph{halts on every input}.

\end{document}